\chapter{结论与展望}
\label{chap:conclusion}

本文研究病例规模在疫情监测中的收益边界。出发点是一个常被默认的判断：报告病例越多，对传播水平的估计越准确。第三章到第六章表明，这一判断对一类估计目标成立，对另一类估计目标只在一定规模以内成立，边界的位置可以折算成病例数，并在住院监测数据中观察到对应的误差形状。本章归纳主要结论，说明研究的局限，并提出监测实践的建议与后续研究方向。

\section{主要结论}
\label{sec:concl-main}

第一，规模的收益取决于估计目标。在负二项传播与二项报告的单代模型下，若目标是当期再生数 $R_t$，估计量 $C_t/(\rho n)$ 是无偏的极大似然估计。它在任意有限规模下的效率都达到最优，方差随规模趋于零（命题~\ref{prop:current}）。若目标是跨期平均再生数 $\Rbar$，方差为 $v_R+A/m$，规模只能消除第二项，环境方差 $v_R$ 构成规模无法突破的底板（命题~\ref{prop:average}）。两类目标使用同一份数据，对规模的需求相差一个数量级以上。

第二，环境底板可以折算成病例数。交叉点 $\mx=A/v_R$ 是规模项等于环境项时的预计报告规模，规模达到 $\mx$ 时误差的一半来自环境波动，扩大规模至多再降低一半误差。对住院监测，$A$ 中的超传播项很小，交叉点主要由 $\Rbar$ 与 $v_R$ 决定；按文献给出的情景参数，交叉点为每期 21.3--49.2 例报告。报告率很低时，扩大覆盖与提高报告率的边际效果几乎相同（命题~\ref{prop:reporting}）。

第三，方差结构给出可操作的设计规则。预算固定时应优先增加期数，正自相关只减慢而不逆转这一收益，有效期数近似则会高估环境项（命题~\ref{prop:allocation}）。考虑每期固定成本后，最优单期规模为交叉点与成本比乘积的平方根量级（推论~\ref{cor:optimal}）。规模较小时，基于负二项分布的精确检验在实际水平 $\alpha^*$ 下一致最优（命题~\ref{prop:test}）。要稳定检出 15\% 的增长，所需规模约为示例交叉点的十倍。更新过程仿真表明，在时间序列中单代公式只能作参照曲线；以实际规模为横轴时，误差形状与参照曲线大体一致。

第四，住院监测数据呈现预期的误差形状。美国 51 个辖区流行初期的州级流感住院周报中，单步预测的相对均方误差随实际规模下降并趋于常数。截距为 0.0991，斜率为 2.243，所有重抽的斜率都为正，经验交叉点为 22.6 例/周。截距与斜率都是情景隐含值的三倍以上，比值的吻合不构成结构参数的检验。分母中的二项报告噪声按滞后权重平方和抬高斜率而不改变截距，能解释斜率水平的一半左右；截距随分母回溯滞后上升，与环境的时间相关和周增长比的口径一致。

第五，规模效应在多种设定下稳健，交叉点的数值则不稳定。工具变量估计、事前基准、时变基准与两组外部数据下，规模项都为正；经验交叉点在 10.6--124.7 之间变化，只能作为量级参考。达峰期以旧基准预测时，88.9\% 的误差来自基准偏移，规模效应被掩盖；自愿填报期周增长比的标准差扩张 192\%，稳健离散度不变，扰动来自少数极端周。全国汇总序列的误差约为州级截距的三分之一，表明空间汇总能有效降低环境底板。

上述结论对应第一章提出的四个子问题。估计目标对规模收益的影响由命题~\ref{prop:current} 与命题~\ref{prop:average} 回答：当期再生数的抽样误差随规模趋于零，跨期平均的误差存在环境底板。规模收益的边界由交叉点 $\mx$ 度量，住院监测的情景交叉点为每期数十例报告。方差结构对监测设计的含义由命题~\ref{prop:reporting}、命题~\ref{prop:allocation}、推论~\ref{cor:optimal} 与精确检验的功效计算给出。单代结论在时间序列中的适用程度，由第四章的更新过程仿真与第五章、第六章的实证共同回答。误差形状在真实监测数据中可以观察到。交叉点的具体数值则受规模测度、基准与病种的影响，尚不能作为精确的设计参数直接使用。

\section{主要创新点}
\label{sec:concl-novelty}

本文的创新之处体现在三个方面。

一是在负二项传播与二项报告的框架下区分当期再生数与跨期平均两个估计目标，并给出两者在有限规模下的精确方差结论。当期目标的有限样本最优性与跨期目标的方差分解并列给出，规模收益的上限由此得到明确表述。

二是把环境底板折算为病例数刻度。交叉点 $\mx=A/v_R$ 把环境方差与规模项放在同一尺度上，期数与规模的预算分配、成本最优单期规模都可以用它表达，后者与两阶段抽样的费用最优分配同构，便于与抽样调查的设计经验相互印证。

三是在真实监测数据上检验误差形状，并分析拟合系数偏离理论值的来源。实证部分用滞后权重的比较核对分母噪声的一阶近似，用达峰期与填报制度变化识别基准偏移与结构性扰动两类偏离，并以工具变量、事前基准、时变基准、空间汇总与外部数据检验结论的稳健性。

\section{研究局限}
\label{sec:concl-limits}

以下七点限定了上述结论的适用范围。

周增长比与单代再生数的口径不同。流感的世代间隔约 3 天，$G_t$ 跨两代左右，其均值与方差都不同于模型中的 $R_t$。实证部分只检验了误差随规模变化的形状，拟合的 $a$、$b$ 不能解读为 $v_R/\Rbar^2$ 与 $A/\Rbar^2$ 的估计。估计单代参数需要日度数据或显式的世代间隔反卷积。

分母含有报告噪声。$D_t$ 是带报告噪声的计数，第五章第~\ref{sec:emp-inflation}~节表明，它抬高斜率而不抬高截距，使 $b/a$ 系统性地高估 $\mx$。二项噪声只解释了实证斜率的一半左右，剩余部分可能来自报告过程的超额离散，也可能来自各州住院比例的差异等其他来源，现有数据无法区分。经验交叉点随滞后权重在 12.6--46.0 之间变化，22.6 只能作为量级。

单代公式在时间序列中不构成下界。第四章式~\eqref{eq:mse-decomp} 的第一项可能高于也可能低于单代公式，更新过程仿真中两种情形都出现了。全文把单代公式作为参照曲线，没有用它判断实证误差是否已经到达极限。

环境存在时间相关。理论部分的假设 H4 要求规模与当期环境无关，更新过程仿真也让各周 $R_t$ 独立。真实的 $R_t$ 有时间相关，大规模往往出现在传播较快的时期。实证截距随回溯滞后上升，第六章的时变基准在大规模处降低了误差，两者都是这种相关的表现。全文没有对 $R_t$ 的相关结构建模，相关程度如何决定截距的大小尚未得到定量回答。

推断方法有两处不足。聚类重抽只在州或周一个维度上进行。州与周的双向聚类在聚类数较少时方差估计不稳定\citep{cameron2011robust, cameron2015practitioner}，样本有 51 个州，而每个流行初期窗口只有 8--11 周，因此只报告了两个单向结果。工具变量回归的排除限制无法检验，其结果只能作为替代性规模测度下的稳健性证据。

部分参数依赖外部取值。情景区间依赖离散参数 $k$ 与住院比例 $\rho$ 的外部取值，二者在流感住院数据中都没有精确估计。精确增长检验要求 $k$ 已知，以估计值代入后检验水平不再有精确保证。由全国序列推算的州间相关 $\bar c\approx0.34$ 忽略了分母噪声与口径问题，只是粗略的量级。

样本窗口是事后划定的。第六章的事前基准只解决了基准的前视问题，流行初期与达峰期窗口的划定仍然用到了全国序列的事后信息。第二期缺失的 36 个州周按 0 处理，在自愿填报期，未填报与零住院不能等同，这影响第二期离散程度的解读。

\section{对监测实践的建议}
\label{sec:concl-practice}

监测问题应先明确估计目标。判断本周疫情是否在增长、及时发出预警，属于当期再生数的估计或检验，规模越大越好，资源应投向报告及时、覆盖完整的单元。病例较少时，宜采用基于负二项分布的精确检验，并按第四章的功效计算评估现有规模能否检出目标增幅。估计一个流行季或一段时期的平均传播水平，属于跨期平均的估计，规模达到交叉点的量级后，继续扩大病例捕获的回报迅速递减。

趋势类问题在规模达到交叉点量级后，应把资源转向增加观测期数与汇总辖区。总报告量固定时，分成更多的期数比集中在少数几期更有效；有每期固定成本时，按推论~\ref{cor:optimal} 选择单期规模。辖区间的环境波动只有部分同步，汇总多个辖区能显著降低环境底板，全国序列的误差约为州级截距的三分之一。区域性的趋势判断宜以汇总序列为主，单个辖区的序列更适合用于预警。

流行阶段转换时应优先更新基准模型。达峰期以旧基准预测，绝大部分误差来自基准偏移，扩大规模无济于事；能跟踪阶段变化的模型，例如滑动窗口估计或集成预测，可以显著降低偏差项。时变模型在小规模时自身噪声较大，规模较小的辖区可以借助汇总或平滑来控制这部分噪声。

数据质量监控宜同时使用标准差与稳健离散度。填报制度变化带来的扰动集中在少数极端周，标准差与偏度对其敏感，四分位距与中位绝对偏差反映整体噪声水平。两类指标的走势出现背离时，提示数据中存在结构性扰动，应核查填报完整性，而不应直接解读为传播水平的变化。

\section{研究展望}
\label{sec:concl-outlook}

后续研究可以沿以下方向展开。

一是使用日度数据与世代间隔反卷积。周增长比混合了约两代的传播，若能获得日度报告数据并显式地按世代间隔分布反卷积，可以直接估计单代参数 $v_R$ 与 $A$，把经验交叉点与理论交叉点放在同一口径下比较。

二是对再生数的时间相关建模。第五章与第六章的证据都指向 $R_t$ 的时间相关。把环境层设为自回归或状态空间过程，并允许规模与环境相关，可以推导截距随回溯滞后变化的定量关系，也能更准确地刻画时变基准降低误差的上限。

三是在精确检验中处理离散参数的不确定性。精确检验目前假定 $k$ 已知。用条件检验、贝叶斯预测分布或对 $k$ 取保守值的方法，可以在 $k$ 需要估计时控制检验水平，并评估功效损失。

四是改进推断方法。随着流行季数增加，州与周的双向聚类会变得可行；也可以采用针对少聚类数的野聚类自助法\citep{cameron2008bootstrap}。工具变量的排除限制可以借助更多的外生变动来源，例如填报规则在不同州的差异，加以检验或替代。

五是用预先登记的窗口做前瞻验证。流行初期与达峰期窗口目前由事后信息划定。在新的流行季开始前按规则预先登记窗口与分析方案，再用实时数据检验误差形状与交叉点，可以排除窗口选择带来的偏差，也能检验交叉点作为监测设计参数的实用价值。

此外，方差分解本身不限于流感住院。门诊流感样病例、实验室病原检测与污水病毒载量等监测数据流，报告机制与离散程度各不相同，规模项 $A$ 的构成也不同。把模型中的二项报告层替换为相应的观测模型，可以为每一类数据流计算交叉点，并据此比较不同数据流在预警与趋势估计两类任务中的相对价值，为多源监测的资源配置提供依据。
